\section{Operator meaning and Route-A decision}

The obstruction in \cref{thm:obstruction} concerns the most natural
unrenormalized symbol. It does not assume that the symbol already comes
from a complete scattering theory. In fact, a second obstruction prevents
that promotion by the ordinary Fredholm route.

\begin{proposition}[Multiplier no-go]\label{prop:multiplier}
Let $F$ be a measurable matrix function on a non-atomic measure space. The
multiplication operator $M_F$ on vector-valued $L^2$ is compact only if
$F=0$ almost everywhere. Hence a nontrivial Mellin multiplier $S_H-I$ is
not trace class.
\end{proposition}

\begin{proof}
If $F$ is nonzero on a set of positive measure, restrict to a subset on which
one fixed matrix entry $|F_{ij}|$ is bounded below. Split that subset into
countably many disjoint positive-measure sets and multiply their normalized
indicators by the fixed basis vector $e_j$. They form an orthonormal sequence,
while the $i$th components of their images have norms bounded below.
Compactness fails.
\end{proof}

Thus the scalar function \eqref{eq:detS} is a fiber determinant, not an
ordinary Fredholm determinant on global Mellin $L^2$. A Birman--Krein
determinant would need a constructed scattering pair and trace-class
resolvent difference. A semifinite or crossed-product determinant would
need a specified algebra and trace. Neither follows from parity
diagonalization.

The strict Route-A classification of the frozen candidate is
\begin{equation}\label{eq:route}
(A1\_\mathrm{WEAK},A2\_\mathrm{FAIL},A3\_\mathrm{FAIL},
A4\_\mathrm{NATURAL\_QUANTIZATION}),
\end{equation}
with overall decision
\begin{center}
\texttt{ROUTE\_A\_REJECTED\_FOR\_UNRENORMALIZED}\\[-2pt]
\texttt{MELLIN\_PARITY\_CANDIDATE}.
\end{center}
The H\'enon phase and its quantization are intrinsic, but there is no
ordinary determinant and the certified extra divisor lies where the target
completed Riemann function is certified nonzero. Route B is not invoked.

The odd parity factor does not rescue the construction. The completed
Riemann zeta function belongs to the trivial real-place parity, which is the
even channel. Keeping only $B$ changes the target archimedean character.
Similarly, dividing by a factor designed from the computed zero would fit
the divisor after the fact. Only a reference operator derived independently
from the same dynamics could authorize cancellation, and the linear parent
\eqref{eq:linear} supplies none.
